% ==========================================================================
% GENERATED by scripts/split_paper.py. Do not edit this file.
%
% Source: imported_papers/grokking/iclr2027_conference_v2.tex
% Part 14: The two endpoint rules in detail
%
% Edits belong upstream, in the published paper. To change what the
% chapter includes, change imports/grokking.toml and re-run
% the splitter.
% ==========================================================================
\section{The two endpoint rules in detail}
\label{app:figs}

\begin{figure}[!ht]
\begin{center}\includegraphics[width=0.62\linewidth]{figures/fig3_steer.pdf}\end{center}
\caption{Turning the knob moves the circuit continuously while nothing else moves. Circuit participation and deletion tolerance against the knob strength $\beta$ on modular addition at $p=31$, three seeds, with test accuracy, logit margin and training loss held across the sweep. The star marks where DFC lands without any knob.}
\label{fig:steer}
\end{figure}

\begin{figure}[!ht]
\begin{center}\includegraphics[width=0.52\linewidth]{figures/fig1_grok.pdf}\end{center}
\caption{Backpropagation and DFC grok to perfect test accuracy after a delay. Train and test accuracy over training for both rules on modular addition, same initialization (faint = train, solid = test).}
\label{fig:grok}
\end{figure}

Backpropagation fixes the reference circuit, since it is the rule every published study of the grokked circuit uses, and DFC is the working rule furthest from it. Both show the grokking signature (Figure~\ref{fig:grok}). Training accuracy saturates by epoch $250$ for both, while test accuracy sits at chance for hundreds of epochs more, crossing $0.95$ at epoch $1500$ for backpropagation and $750$ for DFC on this seed.

Three measures show the difference (Figure~\ref{fig:endpoints}). The concentrated and distributed split appears in the input embedding spectra of panel (a) as well as over units. Panel (b) has participation, spectral entropy and per-neuron periodicity rising together, and periodicity matters because a circuit could spread by making every unit respond to many frequencies, which would lower it. It rises instead, so DFC spreads the circuit while each unit tunes to one frequency more sharply.

Panel (c) shows the gap growing with problem size. The number of frequencies DFC needs for $95\%$ accuracy rises from about $10$ at $p=23$ to $18$ at $p=59$, while backpropagation's count stays flat, slope $0.03$ with interval $[-0.12, 0.18]$ and a slope difference of $0.16$ with $95\%$ interval $[0.04, 0.28]$. That is again a difference in the layout and not in the computation both rules perform (Appendix~\ref{app:algo}). We count frequencies by adding them in order of aggregate power until accuracy reaches $95\%$ of its full value, over four seeds at each modulus.

\begin{figure}[!ht]
\begin{center}\includegraphics[width=\linewidth]{figures/fig2_endpoints.pdf}\end{center}
\caption{The two rules realize the same algorithm over different numbers of units. (a) input embedding Fourier spectra, sparse peaks for backpropagation versus a flat spectrum for DFC. (b) participation ratio, per-neuron periodicity, and spectral entropy. (c) number of frequencies for $95\%$ accuracy versus $p$, where DFC's count grows faster than backpropagation's (slope difference $0.16$, $95\%$ CI $[0.04, 0.28]$).}
\label{fig:endpoints}
\end{figure}

\begin{figure}[!ht]
\begin{center}\includegraphics[width=0.55\linewidth]{figures/fig5_creditflow.pdf}\end{center}
\caption{Both rules concentrate credit during the grokking transition, but backpropagation's concentration is deep and persists while DFC's is shallow and transient. Input layer credit participation over training (four seeds, bands are one standard deviation).}
\label{fig:creditflow}
\end{figure}

Credit over training shows the mechanism in time (Figure~\ref{fig:creditflow}, four seeds). Backpropagation's input layer credit PR collapses from near $0.97$ to about $0.48$ and settles near $0.60$, so the concentration stays in the grokked solution, while DFC dips only to about $0.78$ and returns to $0.93$, passing through the transition and coming back distributed. Concentrating credit is therefore part of how grokking happens under both rules, and what tells the rules apart is whether the concentration is released once the circuit is found.

\begin{figure}[!ht]
\begin{center}\includegraphics[width=\linewidth]{figures/fig4_fragility.pdf}\end{center}
\caption{Distributed circuits keep far more accuracy under unit deletion, while DFC's fuller frequency ensemble is the most robust to frequency removal. Deletion tolerance for backpropagation, the credit knob at $\beta=1$, and DFC. \textbf{(a)} test accuracy as a fraction of hidden units is randomly dropped. \textbf{(b)} test accuracy as input frequencies are randomly removed. \textbf{(c)} area under the damage curve by attack type. Targeted deletion removes the highest norm units first and magnitude pruning removes the largest weights, both deterministic, so the ordering under those attacks cannot be explained by averaging over random masks.}
\label{fig:fragile}
\end{figure}

Damage separates the two ends (Figure~\ref{fig:fragile}, Table~\ref{tab:dissoc}). Under random unit deletion the distributed circuits rank by participation, the knob gaining $+0.07$, $+0.10$ and $+0.14$ over backpropagation at $p = 23, 31, 59$. Under random frequency removal the ranking reverses, DFC gaining $+0.16$, $+0.16$ and $+0.11$, since only DFC spreads the computation over frequencies as well as units, and each spread protects only against the damage that attacks it. DFC's dropout advantage shrinks with $p$ and is gone at $p = 59$, so robustness follows from how the credit spread the circuit, not from which rule trained it.

Distributed circuits also survive targeted deletions better (Table~\ref{tab:targeted}), the knob at $0.58$ and DFC at $0.52$ against backpropagation's $0.33$ when the highest norm units go first, and pruning the largest input weights first leaves both of them above backpropagation as well. A concentrated circuit depends on exactly the units a targeted attack removes first, so no averaging argument explains the gap.

The gap and its consequence hold on all three operations (Table~\ref{tab:ops}). Multiplication is the sharpest case, where backpropagation lands at participation $0.43$ against the knob's $0.99$, so the effect is largest where backpropagation concentrates hardest (Appendix~\ref{app:scope}).

\begin{table}[!ht]
\caption{Under deterministic attacks the distributed rules' advantage widens, so no averaging argument explains it. Deletion tolerance by damage type at $p=31$ (three seeds). Targeted deletion removes the highest outgoing norm units first, and pruning zeroes the largest input weights first. This is an independent run from Table~\ref{tab:dissoc}, so its random deletion row differs from that table's $p=31$ values.}
\label{tab:targeted}
\begin{center}
\small
\begin{tabular}{lccc}
\toprule
Damage & backprop & knob ($\beta{=}1$) & DFC \\
\midrule
Random unit deletion   & $0.674$ & $\mathbf{0.770}$ & $0.725$ \\
Targeted unit deletion & $0.332$ & $\mathbf{0.581}$ & $0.518$ \\
Magnitude pruning      & $0.237$ & $0.358$ & $\mathbf{0.400}$ \\
\bottomrule
\end{tabular}
\end{center}
\end{table}

\begin{table}[!ht]
\caption{The distributed rules keep more accuracy under every damage type and at every modulus. Deletion tolerance by damage type, rule, and modulus (mean and standard deviation over three seeds). Bold marks the most robust rule per row pair. At $p = 23$ backpropagation and the knob do not fully grok within the budget, which inflates the variances there.}
\label{tab:dissoc}
\begin{center}
\small
\begin{tabular}{llccc}
\toprule
Damage & $p$ & backprop & knob ($\beta{=}1$) & DFC \\
\midrule
Unit dropout      & $23$ & $0.576 \pm 0.038$ & $0.648 \pm 0.055$ & $\mathbf{0.674 \pm 0.022}$ \\
                  & $31$ & $0.671 \pm 0.013$ & $\mathbf{0.767 \pm 0.020}$ & $0.716 \pm 0.011$ \\
                  & $59$ & $0.703 \pm 0.013$ & $\mathbf{0.847 \pm 0.002}$ & $0.699 \pm 0.003$ \\
\midrule
Frequency removal & $23$ & $0.532 \pm 0.044$ & $0.513 \pm 0.044$ & $\mathbf{0.696 \pm 0.022}$ \\
                  & $31$ & $0.622 \pm 0.009$ & $0.652 \pm 0.016$ & $\mathbf{0.783 \pm 0.015}$ \\
                  & $59$ & $0.888 \pm 0.013$ & $0.951 \pm 0.007$ & $\mathbf{0.996 \pm 0.001}$ \\
\bottomrule
\end{tabular}
\end{center}
\end{table}

\begin{table}[!ht]
\caption{The participation gap between the rules, and its deletion consequence, replicate on all three operations. Test accuracy, circuit participation, and deletion tolerance for each rule ($p=31$, three seeds, standard deviations in parentheses, training fraction $0.60$ for addition and $0.58$ for subtraction and multiplication).}
\label{tab:ops}
\begin{center}
\begin{tabular}{llccc}
\toprule
Operation & Rule & Test & Circuit PR & Random del \\
\midrule
Addition       & backprop & $0.998\,(0.001)$ & $0.60\,(0.02)$ & $0.671\,(0.016)$ \\
               & knob ($\beta{=}1$) & $0.976\,(0.011)$ & $0.99\,(0.00)$ & $0.771\,(0.019)$ \\
               & DFC & $1.000\,(0.000)$ & $0.92\,(0.00)$ & $0.717\,(0.014)$ \\
\midrule
Subtraction    & backprop & $0.999\,(0.001)$ & $0.59\,(0.04)$ & $0.680\,(0.026)$ \\
               & knob ($\beta{=}1$) & $0.994\,(0.005)$ & $0.99\,(0.00)$ & $0.785\,(0.014)$ \\
               & DFC & $1.000\,(0.000)$ & $0.93\,(0.01)$ & $0.735\,(0.010)$ \\
\midrule
Multiplication & backprop & $0.999\,(0.001)$ & $0.43\,(0.07)$ & $0.666\,(0.035)$ \\
               & knob ($\beta{=}1$) & $0.999\,(0.001)$ & $0.99\,(0.00)$ & $0.867\,(0.012)$ \\
               & DFC & $1.000\,(0.000)$ & $0.86\,(0.02)$ & $0.764\,(0.010)$ \\
\bottomrule
\end{tabular}
\end{center}
\end{table}
